% !PRINTED 424
\weberpage{424}
%<<CONT
\weq{3}{\Sigma N' = 2\sum^{\delta}\sum_{\partial' > a'}^{\partial'} \frac{\partial'}{e'}\varphi(e') + \sum^{n'} \varphi(\sqrt{n'}),}
wenn $\delta^2$ alle quadratischen Faktoren von $n$ (mit etwaiger
Ausnahme von $n$ selbst) durchläuft.

Setzen wir nun
\[
\partial'\delta = \partial, \quad a'\delta = a, \quad e'\delta = e,
\]
so ist $a\partial = n$, und $e$ ist der größte gemeinschaftliche Teiler
von $a$ und $\partial$; zugleich hat $\partial' > a'$ zur Folge, daß
$\partial > a$ ist. Umgekehrt erhält man aus jeder Zerlegung $a\partial$
von $n$ und jedem Teiler $e'$ von $e$ eine Zerlegung $a'\partial'$ von
$n'$, wobei
\[
\delta = \frac{e}{e'}, \quad \partial' = \frac{\partial e'}{e}, \quad a' = \frac{a e'}{e}
\]
wird. Demnach ist
\weq{4}{\Sigma N' = 2\sum_{\partial>a}^{\partial} \frac{\partial}{e} \sum \varphi(e') + \sum^{n'} \varphi(\sqrt{n'}).}

Hierin machen wir nun Gebrauch von der zahlentheoretischen Relation:
$\sum \varphi(\partial) = n$, worin sich die Summe auf alle Divisoren von
$n$ bezieht. Nehmen wir zunächst an, daß $n$ kein Quadrat sei, so sind
auch alle $n'$ keine Quadrate und $\varphi(\sqrt{n'})=0$. Also erhalten
wir wegen $\sum \varphi(e') = e$
\weq{5}{\Sigma N' = 2\sum_{\partial>\sqrt{n}} \partial.}

Wenn dagegen $n$ ein Quadrat ist, so durchläuft $e'$ noch immer die
sämtlichen Divisoren von $e$, $\sqrt{n'}$ aber die sämtlichen Divisoren
von $\sqrt{n}$, mit Ausnahme von 1; danach ergibt sich für diesen Fall:
\weq{6}{\Sigma N' = 2\sum_{\partial>\sqrt{n}} \partial + \sqrt{n} - 1.}

In beiden Fällen durchläuft $\partial$ die sämtlichen Divisoren von $n$,
die \gesperrt{größer} als $\sqrt{n}$ sind.

Wir haben nun noch die Summe der Ausdrücke $\Sigma k h$ für die
verschiedenen Werte von $n'$ oder $\delta$ zu bilden. Setzen wir aber
$n = \delta^2 n'$, so erhalten wir aus jeder eigentlichen Lösung von
\weq{7}{4n' = y'^2 + m x'^2}
durch Multiplikation mit $\delta^2$ eine (wenn $\delta > 1$ ist,
uneigentliche) Lösung von
\weq{8}{4n = y^2 + m x^2,}
nämlich
\[
y = \delta y', \quad x = \delta x'.
\]
